\cotefile{19}
\scanpage{36}
% L1, upper left
\uline{\unsure{XI|XII}} 6.2. \unsure{----|\ldots}

% L2–3
\[
T^i(X)\unsure{\mathord{-}|\to}T'{}^i(X)\unsure{,|;}
\quad \deleted{\text{\unsure{dire pourquoi est} \illegible{1}}}
\quad\begin{aligned}&\text{\unsure{iso} pour }i\le n\\
&\text{mono pour }i=n+1\end{aligned}
\]
\margin{\unsure{!?|?}}

% L4
\[
\unsure{0}\to A-B-C\to 0
\]

% L5–6
\diagram{\begin{tikzcd}[ampersand replacement=\&,column sep=small,row sep=small]
\&T^{i-1}(B)\arrow[r]\arrow[d]\&T^{i-1}(C)\arrow[r,dash]\arrow[d]\&T^i(A)\arrow[r,dash]\arrow[d]\&T^i(B)\arrow[r]\arrow[d]\&T^i(C)\arrow[d]\\
\text{(I)}\&T'{}^{i-1}(B)\arrow[r,dash]\&T'{}^{i-1}(C)\arrow[r,dash]\&T'{}^i(A)\arrow[r,dash]\&T'{}^i(B)\arrow[r,dash]\&T'{}^i(C)
\end{tikzcd}}

% L7–12, a)
a) Supposons qu'on sache que $\varphi^i$ un isom pour $i\le n-1$ \inserted{mono pour $i=n$.} Et qu'il y a un syst. de cog. pour lesquels \deleted{\unsure{$i$}\illegible{1}} $\varphi^n$ est un isom. Alors le diagramme (I) pour $i=n$ montre que \deleted{\unsure{pour}} $\varphi^n$ est un isom pour \uline{tous} les \deleted{\unsure{g}} arguments. Plus \unsure{préc.|gén.}, si $A$ est \inserted{et qu'il $\exists\ A\to B$ mono, et que c'est vrai pour $\varphi^n(B)$, c'est vrai pour $\varphi^n(A)$.}

% L13–17, b)
b) Supposons qu'on sache que $\varphi^i$ est un isom pour $i\le n$, et $\varphi^{n+1}$ un mono pour un syst. de cog. Alors le diagramme (I) avec $i=n+1$ montre que $\varphi^{n+1}$ est un mono. pour \uline{tous} les \uline{arguments}. Plus \unsure{préc.|gén.}, \inserted{si $\exists\ A-B$ mono, et si c'est vrai pour $B$, c'est vrai pour $A$.}

% L18–24, a')
a') $\overline{\mathrm M}$ hypothèse préliminaire sur les $\varphi^i$ que en a) (isom si $i\le n-1$, mono si $i=n$). On veut savoir si $\varphi^n(B)$ est un \uline{isom} (i.e. un \unsure{épim|section}) pour $B$ donné. Si on a une suite exacte $0\to A\to B\to C\to 0$, il suffit de le vérifier pour $\varphi^n(A)$ et $\varphi^n(C)$ ? Oui, \uline{pourvu que} $\varphi^{n+1}(A)$ soit injectif, \uuline{non} en général.

% L25–27, b')
b') $\overline{\mathrm M}$ hyp \unsure{préliminaire} que b) ($\varphi^i$ un isom si $i\le n$). On veut savoir si $\varphi^{n+1}(B)$ est un mono. Suffit-il de le vérifier pour $\varphi^{n+1}(A)$ et $\varphi^{n+1}(C)$ ? \uuline{Oui}.

% archive sheet number: 35
